% Introduction (main text). Numbers belong in the figures and their captions; the text states what they show.
\section{Introduction}
\label{sec:intro}

\IEEEPARstart{S}{ingle-cell} studies now routinely profile hundreds of thousands to millions of cells, reference atlases
hold tens of millions~\cite{svensson2018exponential,sikkema2023integrated,czi2025cellxgene}, and a single perturbation
screen has reached a hundred million~\cite{zhang2025tahoe}. These data are no longer analysed once on a laptop. Atlases
serve them to thousands of users, and machine-learning models stream them minibatch by minibatch, reading every cell many
times~\cite{lopez2018deep,theodoris2023transfer,cui2024scgpt}. More and more often they live in cloud object storage and
cross a network before they are used. Moving the data, rather than computing on them, has become a bottleneck for
large-scale analysis and for building AI-ready datasets.

For a link of a given bandwidth, the number of cells delivered per second is set by how many bits each cell takes. That
number has a floor, the information the counts actually contain, and the formats in use stay far above it.
AnnData~\cite{virshup2024anndata}, the 10x Genomics HDF5 file, loom, Zarr and TileDB-SOMA~\cite{tiledbsoma} store a
matrix of whole-number counts as generic arrays and compress them with tools such as gzip or zstd, which see only bytes.
What limits how fast single-cell data can be delivered today is therefore mostly how they are encoded, not the network.

Count matrices are far from arbitrary. Most non-zero counts are exactly one; how deeply a cell was sequenced sets the
expected count of every gene in it; and if genes are numbered from most to least frequently detected, the positions of a
cell's non-zero entries become a short list of small steps. A code that knows these regularities can describe the same
matrix in far fewer bits, as domain-specific formats for sequencing reads did when they replaced gzip with models of their
data~\cite{bonfield2022cram}. But fewer bits help only if they can be decoded as fast as the link delivers them. Two
limits therefore matter together: how few bits a lossless code can spend on the counts, and how fast a code that small can
be decoded.

We measured both limits on real data and built a format that approaches them. The result is \sone, an open format for
single-cell count matrices, read and written from Python, R and C++ by the \texttt{singlet} package. Fig.~\ref{fig:wow}
summarises the outcome. This paper makes four contributions:
\begin{itemize}
\item \emph{Limits.} We estimate the information floor of \nm{nMatrices} public count matrices with cross-validated count
  models, and the speed frontier of the integer codes a format could adopt, each decoding the same matrices on one core
  (Section~\ref{sec:floor}).
\item \emph{A format.} \sone{} models the counts with a static, context-conditioned entropy coder, stores cells in
  independently decodable blocks behind an index, and names every file by its content (Section~\ref{sec:format}).
\item \emph{A benchmark.} We compare \sone{} with fifteen storage formats and settings on matrices from
  \nm{nProtocolClasses} protocol classes, from cold reads to training loaders on emulated and real cloud storage, with
  every number produced by one open pipeline (Section~\ref{sec:setup}).
\item \emph{Findings.} \sone{} files are the smallest on every matrix and within a few percent of the best count models; no
  measured code was both as small and faster to decode; and from cloud storage \sone{} keeps a network link busy where
  chunked formats leave it idle (Section~\ref{sec:results}). We also report its costs.
\end{itemize}
The codec was developed iteratively with an AI model, Claude (Anthropic), as a coding assistant, with every candidate
change measured against the limits above.

% Figure 1: the headline. Every value is read from generated/fig1*.tsv and generated/fig1a_boxes.tex (written by
% `python -m s1bench report`). One tikzpicture with the axes placed on a grid, so that panels share edges: a and b are
% square and side by side, c (two axes) sits under a and d under b, with the same height.
% \figbox{pos}{q1}{median}{q3}{low whisker}{high whisker}{colour}: one horizontal box of panel a, in axis coordinates.
\newcommand{\figbox}[7]{%
  \draw[#7!80!black, line width=0.5pt] (axis cs:#5,#1) -- (axis cs:#2,#1);
  \draw[#7!80!black, line width=0.5pt] (axis cs:#4,#1) -- (axis cs:#6,#1);
  \draw[#7!80!black, fill=#7, fill opacity=0.18, line width=0.5pt] (axis cs:#2,{#1-0.32}) rectangle (axis cs:#4,{#1+0.32});
  \draw[#7!70!black, line width=1.2pt] (axis cs:#3,{#1-0.32}) -- (axis cs:#3,{#1+0.32});}
% an axis break across a bar that leaves the scale: two slanted strokes on a white gap at (#1, row #2) in axis units
\newcommand{\figbreak}[2]{%
  \fill[white] ($(axis cs:#1,{#2-0.34})+(-1.5pt,0)$) rectangle ($(axis cs:#1,{#2+0.34})+(1.5pt,0)$);
  \draw[ink2, line width=0.5pt] ($(axis cs:#1,{#2+0.34})+(-3pt,0)$) -- ($(axis cs:#1,{#2-0.34})+(0pt,0)$);
  \draw[ink2, line width=0.5pt] ($(axis cs:#1,{#2+0.34})+(0pt,0)$) -- ($(axis cs:#1,{#2-0.34})+(3pt,0)$);
}
\begin{figure*}[tbp]
\centering
\pgfmathsetlengthmacro{\figS}{0.34\linewidth}      % side of the square panels a and b
\pgfmathsetlengthmacro{\figH}{0.52*\figS}          % height of panels c and d
\pgfmathsetlengthmacro{\figG}{1.75cm}              % horizontal gap: the y label and tick labels of b and d
\pgfmathsetlengthmacro{\figV}{1.45cm}              % vertical gap: the x labels of a and b and the titles of c and d
\pgfmathsetlengthmacro{\figCw}{0.46*\figS}         % width of each half of panel c
\begin{tikzpicture}
% ---------------- (a) compression ratio of every matrix against uncompressed CSC ----------------
\begin{axis}[name=pa, at={(0,0)}, anchor=north west, scale only axis, width=\figS, height=\figS,
  title={a\quad Compression of every matrix},
  xmin=0, xmax=17, xtick={0,2,...,16}, xlabel={compression ratio against 32-bit CSC arrays ($\times$)},
  ytick=data, ytick style={draw=none}, yticklabels from table={\gen{fig1a_size.tsv}}{label}, y dir=reverse, enlarge y limits=0.04,
  axis x line*=bottom, axis y line*=left, clip=false]
\addplot[draw=none, forget plot] table[x=ratio_median, y=pos] {\gen{fig1a_size.tsv}};   % sets the y ticks
\draw[gen!60, line width=0.4pt, dashed] (axis cs:1,-0.5) -- (axis cs:1,15.5);
\input{\gen{fig1a_boxes.tex}}
\expandafter\pgfplotsinvokeforeach\expandafter{\otherformats,s1}{
  \addplot[only marks, mark=*, mark size=0.75pt, c-#1, opacity=0.6, discard if not={format}{#1}] table[x=ratio, y=y] {\gen{fig1a_points.tsv}};
}
\end{axis}
% ---------------- (b) size vs speed with R_eff iso-lines (figures/fig1b_panel.tex) ----------------
\def\fbLog{0}\def\fbPts{1}\def\fbTitle{b\quad Smaller and faster, four cores}\def\fbInPicture{1}
\def\fbAxisOpts{name=pb, at={($(pa.north east)+(\figG,0)$)}, anchor=north west, scale only axis, width=\figS, height=\figS}
% Fig. 1b: stored size against cold read rate, from generated/fig1c_tradeoff.tsv (medians and quartiles per format) and
% generated/fig1b_points.tsv (every matrix). Set before \input: \def\fbLog{1} for log axes or {0} for linear ones, and
% \def\fbPts{1} to draw every matrix (with \def\fbPts{0} each median gets its interquartile ranges instead), and the
% panel title in \fbTitle.
\ifnum\fbLog=1
  \def\fbopts{xmode=log, ymode=log, xmin=3, xmax=110, ymin=2, ymax=2000,
    xtick={5,10,20,50,100}, xticklabels={5,10,20,50,100}, ytick={10,100,1000}, yticklabels={10,100,1000},
    xlabel={stored bits per non-zero (log scale)}, ylabel={cold read rate, million non-zeros/s (log scale)}}
  \def\fbiso{25,50,100,200}\def\fbxlo{3}\def\fbtop{1500}\def\fbanchor{anchor}
\else
  \def\fbopts{xmin=0, xmax=70, ymin=0, ymax=260, xtick={0,10,...,70}, ytick={0,50,...,250},
    xlabel={stored bits per non-zero}, ylabel={cold read rate, million non-zeros/s}}
  \def\fbiso{25,50,100}\def\fbxlo{0.2}\def\fbtop{250}\def\fbanchor{anchor_lin}
\fi
% Inside another tikzpicture (Fig. 1), \fbAxisOpts gives the axis position and size and no picture is opened here.
\def\fbaxis#1{\begin{axis}[#1]}
\ifdefined\fbAxisOpts\else\begin{tikzpicture}\def\fbAxisOpts{width=0.47\linewidth, height=7.6cm}\fi
\edef\fbAll{\unexpanded\expandafter{\fbopts},\unexpanded\expandafter{\fbAxisOpts}}
\expandafter\fbaxis\expandafter{\fbAll,
  title={\fbTitle}, clip mode=individual,
  axis x line*=bottom, axis y line*=left,
  lbl/.style={nodes near coords, point meta=explicit symbolic, visualization depends on={value \thisrow{\fbanchor}\as\lblanchor},
    every node near coord/.append style={font=\sffamily\tiny, anchor=\lblanchor, text=ink2, inner sep=1.5pt}},
  iqr/.style={error bars/x dir=both, error bars/y dir=both, error bars/x explicit, error bars/y explicit,
    error bars/error bar style={line width=0.5pt, opacity=0.8}, error bars/error mark=none},
}
% iso-lines of Eq. (1) at B = 1 Gbit/s: R_eff = R  <=>  D = 1/(1/R - b/1000), D and R in million non-zeros/s, b in bits;
% each ends, labelled, at the top of the plot
\expandafter\pgfplotsinvokeforeach\expandafter{\fbiso}{
  \addplot[domain=\fbxlo:{1000*(1/#1-1/\fbtop)}, samples=150, dashed, sone!55, forget plot] {1/(1/#1 - x/1000)}
    node[pos=1, anchor=south west, font=\sffamily\tiny, text=sone!80!black, inner sep=1pt] {#1\,M/s};
}
\node[anchor=north east, font=\sffamily\tiny, text=sone!80!black, align=right, inner sep=1pt] at (rel axis cs:0.99,0.89)
  {dashed: non-zeros per second\\delivered over a 1\,Gbit/s link};
% labels: beside each median on log axes; on linear axes beside, above or below each point, placed by
% generated/fig1d_labels.tex to avoid overlaps
\tikzset{fblabel/.style={font=\sffamily\tiny, text=ink2, inner sep=0.6pt, outer sep=0pt, fill=white, fill opacity=0.8, text opacity=1}}
\ifnum\fbLog=1\pgfplotsset{fblbl/.style={lbl}}\else\pgfplotsset{fblbl/.style={}}\input{\gen{fig1d_labels.tex}}\fi
\ifnum\fbPts=1
\expandafter\pgfplotsinvokeforeach\expandafter{\otherformats,s1}{
  \addplot[only marks, mark=*, mark size=0.7pt, c-#1, opacity=0.4, forget plot, discard if not={format}{#1}] table[x=bits, y=mnnzs] {\gen{fig1b_points.tsv}};
}
\pgfplotsset{fbmed/.style={draw=white, line width=0.4pt}}
\else
\pgfplotsset{fbmed/.style={iqr}}
\fi
\expandafter\pgfplotsinvokeforeach\expandafter{\otherformats}{
\addplot[only marks, fblbl, mark=*, mark size=1.8pt, c-#1, fbmed, discard if not={format}{#1}]
  table[x=bits, y=mnnzs, meta=label, x error plus expr=\thisrow{bits_q3}-\thisrow{bits}, x error minus expr=\thisrow{bits}-\thisrow{bits_q1},
    y error plus expr=\thisrow{mnnzs_q3}-\thisrow{mnnzs}, y error minus expr=\thisrow{mnnzs}-\thisrow{mnnzs_q1}] {\gen{fig1c_tradeoff.tsv}};
}
\addplot[only marks, fblbl, mark=*, mark size=2.4pt, c-s1, fbmed, discard if not={format}{s1}]
  table[x=bits, y=mnnzs, meta=label, x error plus expr=\thisrow{bits_q3}-\thisrow{bits}, x error minus expr=\thisrow{bits}-\thisrow{bits_q1},
    y error plus expr=\thisrow{mnnzs_q3}-\thisrow{mnnzs}, y error minus expr=\thisrow{mnnzs}-\thisrow{mnnzs_q1}] {\gen{fig1c_tradeoff.tsv}};
\end{axis}
\ifdefined\fbInPicture\else\end{tikzpicture}\let\fbAxisOpts\undefined\fi%

% ---------------- (c) time to the first minibatch, then per minibatch ----------------
% Bars fill half of each row (bar width 0.5 in axis units, rows 1 apart), so gaps equal bars. TileDB-SOMA ML leaves both
% scales: its bar is cut by an axis break and labelled inside with its value.
\pgfplotsset{fig1c/.style={xbar, bar shift=0pt, bar width=0.5, ymin=-0.5, ymax=3.5, y dir=reverse,
  scale only axis, width=\figCw, height=\figH, axis x line*=bottom, axis y line*=left, clip=false,
  ytick={0,1,2,3}, ytick style={draw=none}},
  fig1cval/.style={nodes near coords style={font=\sffamily\tiny, anchor=west, text=ink2}},
  fig1cin/.style={nodes near coords style={font=\sffamily\tiny, anchor=east, text=white}}}
\begin{axis}[fig1c, name=pc, at={($(pa.south west)+(0,-\figV)$)}, anchor=north west,
  title={c\quad First minibatch\ldots}, xmin=0, xmax=0.5, xtick={0,0.2,0.4}, xlabel={seconds},
  yticklabels from table={\gen{fig1b_ttfm.tsv}}{label},
  point meta=explicit, nodes near coords={\pgfmathprintnumber[fixed, fixed zerofill, precision=2, assume math mode=true]{\pgfplotspointmeta}\,s}]
\pgfplotsinvokeforeach{s1,zarr3,zarr}{
  \addplot[fill=c-#1, draw=none, fig1cval, discard if not y={fmt}{#1}] table[x=total, y=pos, meta=total] {\gen{fig1b_ttfm.tsv}};
}
\addplot[fill=c-soma, draw=none, fig1cin, discard if not y={fmt}{soma}] table[x expr={0.5}, y=pos, meta=total] {\gen{fig1b_ttfm.tsv}};
\figbreak{0.33}{3}
\end{axis}
\begin{axis}[fig1c, name=pc2, at={($(pc.north east)+(0.08*\figS,0)$)}, anchor=north west,
  title={\ldots then each minibatch}, xmin=0, xmax=400, xtick={0,200,400}, xlabel={milliseconds}, yticklabels={},
  point meta=explicit, nodes near coords={\pgfmathprintnumber[fixed, precision=0, assume math mode=true]{\pgfplotspointmeta}\,ms}]
\pgfplotsinvokeforeach{s1,zarr3,zarr}{
  \addplot[fill=c-#1, draw=none, fig1cval, discard if not y={fmt}{#1}] table[x expr={1000*\thisrow{per_batch}}, y=pos, meta expr={1000*\thisrow{per_batch}}] {\gen{fig1b_ttfm.tsv}};
}
\addplot[fill=c-soma, draw=none, fig1cin, discard if not y={fmt}{soma}] table[x expr={400}, y=pos, meta expr={1000*\thisrow{per_batch}}] {\gen{fig1b_ttfm.tsv}};
\figbreak{260}{3}
\end{axis}
% ---------------- (d) speed-up over one core, formats that decode in parallel ----------------
\begin{axis}[name=pd, at={($(pb.south west)+(0,-\figV)$)}, anchor=north west, scale only axis, width=\figS, height=\figH,
  title={d\quad Parallel decoding},
  xmode=log, log basis x=2, ymode=log, log basis y=2, xmin=0.85, xmax=9.5, ymin=0.9, ymax=9,
  xtick={1,2,4,8}, xticklabels={1,2,4,8}, ytick={1,2,4,8}, yticklabels={1$\times$,2$\times$,4$\times$,8$\times$},
  xlabel={cores}, ylabel={read rate $\div$ one core},
  axis x line*=bottom, axis y line*=left, unbounded coords=jump,
  legend style={font=\sffamily\tiny, draw=none, fill=none, at={(0.02,0.98)}, anchor=north west, legend columns=1, cells={anchor=west}}]
\addplot[gen!70, densely dotted, no marks, forget plot] table[x=threads, y=ideal] {\gen{fig1d_speedup.tsv}}
  node[pos=0.55, sloped, above, font=\sffamily\tiny, text=gen] {ideal};
\addplot[c-s1, very thick, mark=*, mark size=1.6pt] table[x=threads, y=s1] {\gen{fig1d_speedup.tsv}};
\addplot[c-soma, thick, mark=triangle*, mark size=1.5pt] table[x=threads, y=soma] {\gen{fig1d_speedup.tsv}};
\addplot[c-mtx_gz, thick, mark=square*, mark size=1.2pt] table[x=threads, y=mtx_gz] {\gen{fig1d_speedup.tsv}};
\addplot[c-parquet_coo, thick, mark=diamond*, mark size=1.5pt] table[x=threads, y=parquet_coo] {\gen{fig1d_speedup.tsv}};
\addplot[c-zarr3, thick, mark=pentagon*, mark size=1.4pt] table[x=threads, y=zarr3] {\gen{fig1d_speedup.tsv}};
\legend{\sone, TileDB-SOMA, MTX.gz, Parquet COO, {Zarr v3, sharded}}
\end{axis}
\end{tikzpicture}
\caption{\textbf{\sone{} files are the smallest on every matrix and, on four cores, the fastest to read on
\nm{dec-s1-fastest-count} of \nm{dec-paired-n}.}
\textbf{(a)} Compression ratio of each format's file (counts, gene identifiers and barcodes) against uncompressed CSC
arrays with 32-bit indices and values, one point per matrix (\nm{nMatrices} matrices); boxes, quartiles and median;
whiskers, the most extreme matrices within 1.5 interquartile ranges. \sone{} compresses a median of \nm{cr-s1}$\times$
(\nm{cr-s1-lo}--\nm{cr-s1-hi}$\times$), gzip-compressed AnnData \nm{cr-h5ad_gzip}$\times$.
\textbf{(b)} Size against read rate from a cold file into an in-memory sparse matrix on \nm{threadsHeadline} cores: one
faint point per matrix and format, large points the medians. Dashed lines join sizes and decode rates that deliver the same
number of non-zeros per second over a 1\,Gbit/s link (Eq.~\ref{eq:reff}); up and to the left is better. On one core
\sone{} reads \nm{dec-s1-t1} million non-zeros per second, behind uncompressed arrays (Table~\ref{tab:main}). Log axes:
Supplementary Fig.~\ref{S-fig:sfig1b}.
\textbf{(c)} Seconds from a new process to the first 512-cell training minibatch, decoded into a PyTorch sparse tensor,
and milliseconds per minibatch after that; \nm{fig1b-dataset} read from an object store over a \nm{fig1b-tier} link,
median of three seeds. Bars cut by an axis break carry their value inside. Loaders that download the whole file first are
in Supplementary Table~\ref{S-tab:minibatch}.
\textbf{(d)} Read rate from a cold file relative to one core, median of per-matrix ratios, for the \nm{par-formats-n} of
\nm{par-formats-of} formats whose readers use more than one core.}
\label{fig:wow}
\end{figure*}

